\section{An exact audit of the 2024 rational cylinder construction}
\label{sec:wb-tao2024}

\subsection{Scope and primary-source comparison}
The primary source is T.~Tao, \emph{A possible approach to finite time
blowup for the Euler equations}, 8 September 2024,
\url{https://terrytao.wordpress.com/2024/09/08/a-possible-approach-to-finite-time-blowup-for-the-euler-equations/}.
The retrieved display contains the axial factor $(1+x^2)^{-1}$,
the transverse factor $(1+x^2)^{-2}$, and a middle coefficient
$-2\varphi v$. The calculations below retain these original rational
functions and prove that the coefficient they generate is $-4\varphi v$.
The source's equation (7) also has its two Gram-matrix rows interchanged.
These are statements about the retrieved formulas, not claims about a
singularity theorem. All further formulas in this section are derived here.
The incoming report's equations (16.4)--(16.9), on physical PDF pages
30--31 (printed pages 25--26), have the correct derivative, coefficient,
and Gram order. Its local metric argument is extended below to precise
global and volume statements.

\subsection{The full product equations and the pressure map}
Let $N$ be a smooth compact manifold without boundary with a smooth
Riemannian metric $h$. A positive volume density suffices when $N$ is
nonorientable; when a volume form is used, fix an orientation.
On $M=\mathbb R_x\times N$ retain the product metric
$g=dx^2+h$ and write $u=a\partial_x+U$ and $f=b\partial_x+F$.
Here $a,p,b$ are real-valued smooth functions of $(t,x,y)$ and $U,F$
are smooth sections of the pullback of $TN$. The sign convention is
\begin{equation}\label{eq:wb-euler-sign}
 \partial_tu+\nabla^g_u u=-\operatorname{grad}_g p+f,
 \qquad \operatorname{div}_g u=0.
\end{equation}
In a coordinate chart $y^i$ on $N$, the only nonzero Christoffel
symbols of $g$ are $\Gamma^i_{jk}(h)$: this follows directly from
$g_{xx}=1$, $g_{xi}=0$, $\partial_xh_{ij}=0$ in
$\Gamma^A_{BC}=\frac12g^{AD}(\partial_Bg_{CD}+\partial_Cg_{BD}
-\partial_Dg_{BC})$. In particular
\begin{align}
 \partial_ta+a\partial_xa+U(a)&=-\partial_xp+b,
 \label{eq:wb-axial}\\
 \partial_tU+a\partial_xU+\nabla^h_U U&=-\operatorname{grad}_h p+F,
 \label{eq:wb-transverse}\\
 \partial_xa+\operatorname{div}_h U&=0.
 \label{eq:wb-div}
\end{align}
The last identity also follows from
$\operatorname{div}_g u=(\sqrt{\det h})^{-1}
[\partial_x(\sqrt{\det h}\,a)+\partial_{y^i}(\sqrt{\det h}\,U^i)]$.

There is an exact residual construction on this cylinder. For every
smooth pair $(a,U)$ obeying \eqref{eq:wb-div}, set
\begin{align}
 A(t,x,y)&=\partial_ta+a\partial_xa+U(a),\\
 R(t,x,y)&=\partial_tU+a\partial_xU+\nabla^h_U U,\\
 p(t,x,y)&=q(t,y)-\int_0^x A(t,s,y)\,ds,
 \label{eq:wb-pressure-map}\\
 b&=0,\qquad F=\operatorname{grad}_h p+R,
 \label{eq:wb-force-map}
\end{align}
where $q$ is any smooth real function on the time interval times $N$.
The integral is over a finite interval for every $x$, so differentiation
under it proves smoothness on the open spacetime domain. Its $x$
derivative is $-A$ by the fundamental theorem of calculus. Substitution
in \eqref{eq:wb-axial} and \eqref{eq:wb-transverse} proves every component
of \eqref{eq:wb-euler-sign}. Conversely, a solution with $b=0$ has
$p$ of \eqref{eq:wb-pressure-map}, with $q=p(t,0,y)$, and has the
force \eqref{eq:wb-force-map}. Thus the pressure map describes exactly,
and not merely formally, all such force/pressure pairs for a specified
divergence-free velocity. Changing $q$ generally changes the transverse
force; it is not the usual pressure gauge with the force held fixed.

The cancelled-pressure subclass is $R=0$ and
$F=\operatorname{grad}_h p$. Its closed velocity system is therefore
\begin{equation}\label{eq:wb-cylinder-closed}
 \partial_tU+a\partial_xU+\nabla^h_U U=0,
 \qquad \partial_xa+\operatorname{div}_hU=0.
\end{equation}
On an axial circle, the same construction requires
$\int_0^L A(t,s,y)\,ds=0$ for pressure periodicity. The integral formula
on $\mathbb R$ supplies no automatic periodic realization.

\subsection{Energy and the exact constant-parameter scaling}
Set $e=|U|_h^2/2$. Metric compatibility gives
$(\partial_t+a\partial_x+U)e=\langle R,U\rangle_h$.
Since \eqref{eq:wb-div} holds,
\begin{equation}\label{eq:wb-local-energy}
 \partial_te+\partial_x(ae)+\operatorname{div}_h(eU)
 =\langle R,U\rangle_h.
\end{equation}
Indeed, the two product-rule terms added to the transport derivative
are $e(\partial_xa+\operatorname{div}_hU)=0$. Integration on
$[-L,L]\times N$ yields the exact identity
\[
 \frac{d}{dt}\int_{-L}^L\!\int_N e\,dV_h\,dx
 =-\int_N[ae]_{x=-L}^{x=L}\,dV_h
 +\int_{-L}^L\!\int_N\langle R,U\rangle_h\,dV_h\,dx.
\]
Here the integral of the $N$ divergence vanishes: in an oriented
chart $\operatorname{div}_h(eU)dV_h=d\iota_{eU}dV_h$, and Stokes'
theorem on the compact boundaryless manifold gives zero; a partition
of unity gives the identical density statement. Consequently the
transverse energy is conserved for $R=0$ when the energy is finite,
the axial boundary flux tends to zero, and differentiation and the
$L\to\infty$ limit are justified, for example by uniform integrable
majorants for $e$ and $\partial_te$ on each compact time interval.
Finite energy alone is not substituted for this flux/limit requirement.

For positive constants $\lambda,\mu$, let
\[
 \widetilde a(t,x,y)=\frac\lambda\mu
 a(t/\mu,x/\lambda,y),\qquad
 \widetilde U(t,x,y)=\frac1\mu U(t/\mu,x/\lambda,y).
\]
All three terms of $R$ acquire the factor $\mu^{-2}$; the divergence
acquires the factor $\mu^{-1}$. This follows from the ordinary chain
rule and the $x$ independence of $h$. The axial acceleration acquires
$\lambda/\mu^2$. Hence in the cancelled-pressure subclass
$\widetilde p=(\lambda^2/\mu^2)p(t/\mu,x/\lambda,y)$ and
$\widetilde F=\operatorname{grad}_h\widetilde p$ give the exact
full forced equations as well. This is a symmetry of that subclass
with its force transformed, not an assertion about an arbitrary fixed
force in \eqref{eq:wb-euler-sign}. The energy transformation is exactly
\begin{equation}\label{eq:wb-energy-scale}
 \widetilde E_N(t)=\frac\lambda{\mu^2}E_N(t/\mu),\qquad
 \widetilde E_x(t)=\frac{\lambda^3}{\mu^2}E_x(t/\mu),
\end{equation}
by $x=\lambda\xi$. Equal nonzero transverse energies give
$\mu^2=\lambda$; comparable energies give corresponding two-sided
bounds on $\lambda/\mu^2$. No original dimensional constant or metric
has been changed by these statements.

\subsection{All rational coefficient equations, with the discrepancy retained}
For smooth $\varphi\colon N\to\mathbb R$ and smooth vector fields
$v,w\in\Gamma(TN)$, write
\begin{equation}\label{eq:wb-rational}
 D=1+x^2,\qquad a=\frac\varphi D,\qquad U=\frac{v+xw}{D^2}.
\end{equation}
Differentiation, without altering either denominator, gives
\begin{align}
 \partial_xa&=-\frac{2x\varphi}{D^2},\\
 \partial_xU&=\frac{wD-4x(v+xw)}{D^3}
 =\frac{w-4xv-3x^2w}{D^3},\label{eq:wb-U-derivative}\\
 \nabla^h_U U&=\frac{\nabla_vv+x(\nabla_vw+\nabla_wv)
 +x^2\nabla_ww}{D^4}.
\end{align}
The last formula uses $v(D)=w(D)=0$, including all connection terms
in each $\nabla$. Define the three actual vector fields
\begin{equation}\label{eq:wb-C}
 C_0=\nabla_vv+\varphi w,\quad
 C_1=\nabla_vw+\nabla_wv-4\varphi v,\quad
 C_2=\nabla_ww-3\varphi w.
\end{equation}
The full residual and divergence, for arbitrary $\varphi,v,w$, are
\begin{equation}\label{eq:wb-rational-residual}
 R=\frac{C_0+xC_1+x^2C_2}{D^4},\qquad
 \operatorname{div}_g u=
 \frac{\operatorname{div}_hv+x(\operatorname{div}_hw-2\varphi)}{D^2}.
\end{equation}
Since $D>0$ for every real $x$, a vanishing vector polynomial has all
coefficients zero in any chart. Thus the stationary closed system is
equivalent in both directions to the five equations
\begin{equation}\label{eq:wb-correct-five}
 C_0=C_1=C_2=0,\qquad
 \operatorname{div}_hv=0,\qquad\operatorname{div}_hw=2\varphi.
\end{equation}
There are no omitted derivatives of $\varphi$ in its transverse
equation: the advective derivative there differentiates $U$, not $a$.
Those derivatives occur in the axial pressure equation below.

For the primary display's middle expression set instead
$B_1=\nabla_vw+\nabla_wv-2\varphi v$. The exact relation is
$C_1=B_1-2\varphi v$. In particular, retaining all four other
equations and $B_1=0$ gives
\begin{equation}\label{eq:wb-blog-residual}
 R=-\frac{2x\varphi v}{(1+x^2)^4}.
\end{equation}
The discrepancy disappears exactly when the vector field
$\varphi v$ vanishes identically. The residual construction
\eqref{eq:wb-pressure-map}--\eqref{eq:wb-force-map} maps any such
printed-system solution to full forced Euler with
$F=\operatorname{grad}_h p-2x\varphi v/(1+x^2)^4$.
This proves the precise bridge after the coefficient failure; it makes
no assertion that a nondegenerate solution of either system has been found.
More generally \eqref{eq:wb-rational-residual} gives that bridge for
every triple satisfying only its two displayed divergence equations.

\subsection{Explicit pressure, force tails, and energy of the original ansatz}
For \eqref{eq:wb-rational}, the axial acceleration is
\[
 a\partial_xa+U(a)
 =\frac{-2x\varphi^2+v(\varphi)+xw(\varphi)}{D^3}.
\]
Define
\begin{align}
 A_0&=2\varphi^2-w(\varphi),\qquad B_0=v(\varphi),\\
 J(x)&=\frac{x}{4D^2}+\frac{3x}{8D}+\frac38\arctan x,\\
 p(x,y)&=-\frac{A_0(y)}{4D^2}-B_0(y)J(x).
 \label{eq:wb-explicit-pressure}
\end{align}
The derivative of the three summands in $J$ is respectively
$1/(4D^2)-x^2/D^3$, $3/(8D)-3x^2/(4D^2)$, and $3/(8D)$;
putting these over $D^3$ gives $J'=D^{-3}$.
Also $(D^{-2})'=-4xD^{-3}$. Consequently
$p_x=(xA_0-B_0)/D^3$, the negative of the displayed acceleration.
Thus $p$ and
\begin{equation}\label{eq:wb-explicit-force}
 F=-\frac{\operatorname{grad}_h A_0}{4D^2}
   -J(x)\operatorname{grad}_h B_0
   +\frac{C_0+xC_1+x^2C_2}{D^4},\qquad b=0
\end{equation}
are a completely explicit, smooth stationary force/pressure pair for
every divergence-compatible triple. Every spatial covariant derivative
of this pair is bounded on $\mathbb R\times N$: derivatives in $y$
of its coefficients are bounded by compactness, $J$ is bounded,
and all positive-order $x$ derivatives of $J$ and of the displayed
rational factors are bounded. These are stationary fields, so positive
time derivatives vanish. This proves smooth bounded forcing for this
stationary map; it does not imply bounds for a singularly modulated map.

Its force tails have an exact further restriction. Since
$J(+\infty)=3\pi/16$ and $J(-\infty)=-3\pi/16$,
\begin{equation}\label{eq:wb-force-limits}
 \lim_{x\to\pm\infty}F(x,\cdot)
 =\mp\frac{3\pi}{16}\operatorname{grad}_h B_0.
\end{equation}
The other summands tend to zero uniformly on $N$. If
$\operatorname{grad}_hB_0$ is nonzero at a point, it is bounded away
from zero on a nonempty open set; the uniform limit gives a positive
lower bound for $|F|$ on that set for all sufficiently large $|x|$.
Then $F\notin L^2(M)$. If $\operatorname{grad}_hB_0=0$, the remaining
terms are $O(|x|^{-4})$ or faster, and $F\in L^2(M)$.
For the divergence-compatible triples under consideration,
$\int_{N_j}B_0\,dV_h=\int_{N_j}\operatorname{div}_h(\varphi v)\,dV_h=0$
on every connected component $N_j$. A function with zero gradient is
constant on each component, so here the exact $L^2$ criterion is
$B_0=v(\varphi)=0$.
Adding a different $q(y)$ to the pressure adds $\operatorname{grad}_h q$
to both limits in \eqref{eq:wb-force-limits}; both can vanish only
when $\operatorname{grad}_hB_0=0$ and $\operatorname{grad}_hq=0$.
Thus this tail issue cannot be removed by that pressure choice.

The velocity always has finite total energy. The needed integrals are
\begin{equation}\label{eq:wb-integrals}
 \int_{\mathbb R}\frac{dx}{D^2}=\frac\pi2,\quad
 \int_{\mathbb R}\frac{dx}{D^4}=\frac{5\pi}{16},\quad
 \int_{\mathbb R}\frac{x\,dx}{D^4}=0,\quad
 \int_{\mathbb R}\frac{x^2\,dx}{D^4}=\frac\pi{16}.
\end{equation}
For completeness put $x=\tan s$, $-\pi/2<s<\pi/2$.
The first, second, and fourth integrands become respectively
$\cos^2s$, $\cos^6s$, and $\sin^2s\cos^4s$.
Integration by parts of $\int_{-\pi/2}^{\pi/2}\cos^{2m}s\,ds$
gives $I_m=(2m-1)I_{m-1}/(2m)$ with $I_0=\pi$;
hence $I_1=\pi/2$, $I_2=3\pi/8$, $I_3=5\pi/16$ and
the fourth integral is $I_2-I_3=\pi/16$. Oddness gives the third.
Expanding $|v+xw|^2=|v|^2+2x\langle v,w\rangle+x^2|w|^2$
therefore proves
\begin{equation}\label{eq:wb-exact-energy}
 E_N=\frac\pi{32}\int_N(5|v|_h^2+|w|_h^2)\,dV_h,
 \qquad E_x=\frac\pi4\int_N\varphi^2\,dV_h.
\end{equation}
Algebraic decay, compact transverse geometry, and smooth bounded force
are the properties proved here; rapid decay in Euclidean three-space
is not a consequence of these formulas.

\subsection{Exact one-form and Lie-derivative equations}
Let $\theta=h(v,\cdot)$ and $\eta=h(w,\cdot)$; the source calls the
second one-form $\lambda$. The notation $\eta$ here prevents collision
with the positive axial scaling parameter. For any vector field $X$,
metric compatibility and vanishing torsion give
\begin{align*}
 (\iota_vd\theta)(X)
 &=v[\theta(X)]-X[\theta(v)]-\theta([v,X])\\
 &=h(\nabla_vv,X)-h(\nabla_Xv,v),\\
 \tfrac12d[\theta(v)](X)&=h(\nabla_Xv,v).
\end{align*}
Adding proves
\begin{equation}\label{eq:wb-accel}
 (\nabla_vv)^\flat=\iota_vd\theta+\tfrac12d[\theta(v)].
\end{equation}
Apply this identity to $v+w$, expand every bilinear term, and subtract
the individual $v$ and $w$ identities. The resulting mixed identity is
\begin{equation}\label{eq:wb-polarized}
 (\nabla_vw+\nabla_wv)^\flat
 =\iota_vd\eta+\iota_wd\theta
 +\tfrac12d[\eta(v)+\theta(w)].
\end{equation}
Since $h$ is positive definite, lowering indices is a vector-bundle
isomorphism. The three equations $C_i=0$ are consequently equivalent
to
\begin{align}
 \varphi\eta+\iota_vd\theta+\tfrac12d[\theta(v)]&=0,
 \label{eq:wb-form-zero}\\
 -4\varphi\theta+\iota_vd\eta+\iota_wd\theta
 +\tfrac12d[\eta(v)+\theta(w)]&=0,
 \label{eq:wb-form-one}\\
 -3\varphi\eta+\iota_wd\eta+\tfrac12d[\eta(w)]&=0.
 \label{eq:wb-form-two}
\end{align}
No antisymmetric covariant term has been dropped: the mixed expression
is the symmetric sum of the two accelerations, exactly as in $C_1$.

Define the Lie derivative on one-forms by
$(\mathcal L_v\alpha)(X)=v[\alpha(X)]-\alpha([v,X])$.
Expanding $d\alpha(v,X)$ shows directly that
$\mathcal L_v\alpha=\iota_vd\alpha+d[\alpha(v)]$.
Thus the identical system in Lie-derivative notation is
\begin{align}
 \mathcal L_v\theta-\tfrac12d[\theta(v)]+\varphi\eta&=0,\\
 \mathcal L_v\eta+\mathcal L_w\theta
 -\tfrac12d[\eta(v)+\theta(w)]-4\varphi\theta&=0,\\
 \mathcal L_w\eta-\tfrac12d[\eta(w)]-3\varphi\eta&=0.
 \label{eq:wb-lie-system}
\end{align}
With a specified smooth positive density $\omega$, define divergence
by $\mathcal L_v\omega=(\operatorname{div}_\omega v)\omega$.
The divergence constraints eliminate $\varphi$ exactly by
$\varphi=\frac12\operatorname{div}_\omega w$ and
$\operatorname{div}_\omega v=0$. In the middle one-form equation
the resulting coefficient is therefore
$-2(\operatorname{div}_\omega w)\theta$, rather than
$-(\operatorname{div}_\omega w)\theta$. An arbitrary choice of
$\omega$ requires the precise metric/volume realization proved in
the following section; it is not implicitly equal to $dV_h$.

For completeness the full time-dependent transverse equation also
has an exact metric-free acceleration identity. Put
$\zeta=h(U,\cdot)$, a time- and $x$-dependent one-form on $N$.
The same calculation, with $h$ fixed in $t,x$, gives
\[
 R^\flat=\partial_t\zeta+a\partial_x\zeta
 +\mathcal L_U\zeta-\tfrac12d_N[\zeta(U)].
\]
This equation includes all dynamics of the closed cylinder system;
there is no derivative of $a$ in this transverse one-form identity.

\subsection{Time-dependent concentration gives an explicit residual}
Let $(a_*,U_*)$ be any fixed smooth stationary divergence-free pair
on this same cylinder and put
$R_*=a_*\partial_\xi U_*+\nabla_{U_*}U_*$.
For arbitrary smooth positive functions $L(t),M(t)$, define
\[
 \xi=x/L(t),\qquad
 a_{L,M}=\frac{L(t)}{M(t)}a_*(\xi,y),\qquad
 U_{L,M}=\frac1{M(t)}U_*(\xi,y).
\]
Its divergence is $M^{-1}(\partial_\xi a_*+
\operatorname{div}_hU_*)=0$. Direct differentiation proves the
complete transverse residual
\begin{equation}\label{eq:wb-modulation}
 R_{L,M}=\frac1{M^2}R_*(\xi,y)
 -\frac{M'}{M^2}U_*(\xi,y)
 -\frac{L'}{ML}\xi\partial_\xi U_*(\xi,y).
\end{equation}
Indeed $\xi_t=-L'\xi/L$, the two transport terms have factor $M^{-2}$,
and differentiating $M^{-1}U_*(\xi,y)$ gives exactly the last two
terms. Its axial acceleration, with none of the modulation terms
discarded, is
\begin{align}
 A_{L,M}={}&\left(\frac{L'}M-\frac{LM'}{M^2}\right)a_*
 -\frac{L'}M\xi\partial_\xi a_*\notag\\
 &+\frac L{M^2}\bigl(a_*\partial_\xi a_*+U_*(a_*)\bigr).
 \label{eq:wb-modulation-axial}
\end{align}
Formula \eqref{eq:wb-pressure-map} with this $A_{L,M}$ and
\eqref{eq:wb-force-map} with \eqref{eq:wb-modulation} supply an
exact forced Euler solution on every open time interval where
$L,M>0$ are smooth. No assertion of force regularity at an endpoint
where these functions vanish follows from this construction.

If $R_*=0$ and the transverse-energy scaling factor is held fixed by
$M=\kappa\sqrt L$ with a fixed $\kappa>0$, then
\begin{equation}\label{eq:wb-energy-modulation-defect}
 R_{L,M}=-\frac{L'}{ML}
 \left(\tfrac12U_*+\xi\partial_\xi U_*\right).
\end{equation}
For a nonzero smooth $U_*$ on the entire cylinder, the expression
in parentheses cannot vanish identically. To prove this, on each
half-line the equation is
$\partial_\xi(|\xi|^{1/2}U_*)=0$ in each local bundle component.
Thus $U_*=c_+(y)\xi^{-1/2}$ for $\xi>0$ and
$U_*=c_-(y)(-\xi)^{-1/2}$ for $\xi<0$. Continuity at $\xi=0$
forces both coefficient fields to vanish; the equation at zero also
gives $U_*(0,y)=0$. Hence $U_*$ would be zero everywhere.
Therefore nonconstant $L$ in this pure scaling construction produces
a nonzero residual at every time with $L'\ne0$. This proves only
the defect of pure modulation of a fixed smooth profile. The full
residual and pressure formulas above remain an exact force bridge
and do not exclude more general profile dynamics.

There is a concrete nonzero stationary member, including its
degeneracy. On a flat torus $N=\mathbb R^m/\Lambda$ with its
specified constant positive metric and lattice $\Lambda$, choose
a nonzero constant vector $V$ and take
$\varphi=0$, $w=0$, $v=V$.
All five equations \eqref{eq:wb-correct-five} hold because constant
vector fields have zero covariant derivatives and divergence.
Thus $a=0$, $U=V/(1+x^2)^2$, $p=0$, $f=0$ is a smooth unforced
stationary Euler solution on this product cylinder. Its energy is
$5\pi\operatorname{Vol}_h(N)|V|_h^2/32$ by
\eqref{eq:wb-exact-energy}. The pair $(v,w)$ has rank one, so this
example does not realize the independent-pair Gram condition or a
singularity. It verifies directly that the rational construction has
actual solutions while identifying the exact rank of this example.
